%!TEX TS-program = xelatex
%!TEX encoding = UTF-8
% ==============================================================================
%  可切换配色 · 书籍出版物模板（明日方舟·孤星 · Lonetrail Book Template）
%
%  演示一本完整出版物的骨架：封面 → 扉页 → 版权页 → 目录 → 正文 → 封底。
%  目录：styles/ 版式与家具 | chapters/ 正文 | palettes/ 配色 | tools/ 编译换肤
%
%  换肤只改 documentclass 的 theme=NNN，或运行：
%      python tools/lonetrail.py theme 042
% ==============================================================================
\documentclass[
  trim=185x260,
  theme=000,
  bg=grid,
  fontsize=10pt
]{lonetrailbook}

% ============================================================================
% 封面 / 扉页 / 版权页 / 封底 信息（只改这里）
% ============================================================================
\title{测度论讲义}
\subtitle{Lecture Notes on Measure Theory}
\seriesname{明日方舟 · 孤星学术丛书}
\author{莱茵实验室}
\publisher{Rhine Lab Pioneer Project}
\isbn{978-7-0000-0000-0}
\edition{2026 年 10 月第 1 版}
\volumeline{第一卷 · 可测空间}
\volumelineEN{Volume I · Measurable Spaces}
\coverline{$\displaystyle\int_X f\,\mathrm{d}\mu$}
\watermark{$\displaystyle\int$}
\blurb{%
  本书以可测空间为起点，逐步建立外测度、Carathéodory 扩张与 Lebesgue 积分，
  再进入乘积空间、符号测度与 $L^p$ 空间.全书定理、定义、例题六色编码，
  读者扫一眼即可分辨环境类型；配色可一键切换，正文与封面同步换肤.}

% 页眉丛书行 / 页脚注脚行（留空即不显示）
\renewcommand{\LoneSeriesMark}{孤星学术丛书 · 测度论讲义}
\renewcommand{\LoneFootnoteLine}{\kaishu 明日方舟 · 孤星 LaTeX 书籍出版物模板}

% ============================================================================
% 章扉页槽位：按章号给数学母题、篇名与题记（可覆写）
% ============================================================================
\renewcommand{\LonePlateSymbol}[1]{%
  \ifcase#1\relax $\cdot$%
  \or $\mathcal{A}$% 1 可测空间
  \or $\mu$%        2 外测度
  \else $\mu$%
  \fi}
\renewcommand{\LonePlatePartCN}[1]{%
  \ifcase#1\relax 预备%
  \or 可测空间%
  \or 测度论基础%
  \else 进阶%
  \fi}
\renewcommand{\LonePlateQuote}[1]{%
  \ifcase#1\relax 数学是解释自然现象最可靠的语言.%
  \or 集合论的全部困难，在于如何处理无穷.%
  \or 测度论不过是为「大小」找到一个自洽的定义.%
  \else 严格性是数学的良心.%
  \fi}
\renewcommand{\LonePlateQuoteBy}[1]{%
  \ifcase#1\relax 欧拉%
  \or 康托尔%
  \or 勒贝格%
  \else 外尔%
  \fi}

% 卷隔页槽位（多卷合订本用；本 demo 演示一次）
\renewcommand{\LoneVolumeTitle}[1]{第一卷 · 可测空间}
\renewcommand{\LoneVolumeBlurb}[1]{%
  本书以可测空间为起点，逐步建立外测度、Carathéodory 扩张与 Lebesgue 积分.}

% ============================================================================
\begin{document}

\maketitle
\frontmatter

\tableofcontents
\clearpage
{\thispagestyle{empty}
\begin{center}\vspace*{\fill}
{\kaishu\fontsize{13}{22}\selectfont\color{InkMuted}
  本书为书稿版本，内容与版式可能在作者后续修订中调整.\\[1em]
  未经书面许可，不得以任何形式复制、转载或倒卖.}
\vspace*{\fill}\end{center}
\clearpage}

\mainmatter

\VolumeDivider{1}

% 可测空间
\section{$\sigma$-代数}

\begin{wenxintishi}[本章在做什么]
集合论给出「哪些子集值得测量」，$\sigma$-代数把这件事变成一条封闭性公理.
先承认一部分子集可测，再要求它们在补、可数并、可数交下仍然可测，就得到一个
足够大又能自洽运作的舞台.
\end{wenxintishi}

\begin{definition}[$\sigma$-代数]
设 $X$ 为非空集合，$\mathcal{A}\subseteq\mathcal{P}(X)$. 若 $\mathcal{A}$ 满足
\begin{enumerate}[label=\roman*.]
  \item $X\in\mathcal{A}$；
  \item $A\in\mathcal{A}\Rightarrow X\setminus A\in\mathcal{A}$；
  \item $A_1,A_2,\dots\in\mathcal{A}\Rightarrow \bigcup_{n=1}^{\infty}A_n\in\mathcal{A}$，
\end{enumerate}
则称 $\mathcal{A}$ 为 $X$ 上的一个 \emph{$\sigma$-代数}，称 $(X,\mathcal{A})$ 为可测空间，
称 $\mathcal{A}$ 中的元素为可测集.
\end{definition}

\begin{theorem}[$\sigma$-代数的最小生成]
任一非空集族 $\mathcal{E}$ 都被某个 $\sigma$-代数包含，且所有包含 $\mathcal{E}$ 的
$\sigma$-代数之交仍为包含 $\mathcal{E}$ 的 $\sigma$-代数. 记之为
$\sigma(\mathcal{E})$，称为由 $\mathcal{E}$ 生成的 $\sigma$-代数.
\end{theorem}

\begin{proof}
记
\[
\mathcal{F}=\bigcap\{\mathcal{A}\supseteq\mathcal{E}:\mathcal{A}\text{ 是 }\sigma\text{-代数}\}.
\]
由 $\mathcal{P}(X)$ 本身是 $\sigma$-代数，交非空. $\sigma$-代数族对任意交封闭，故
$\mathcal{F}$ 是 $\sigma$-代数，且显然 $\mathcal{E}\subseteq\mathcal{F}$. 由构造它是最小的.
\end{proof}

\section{可测函数}

\begin{definition}[可测映射]
设 $(X,\mathcal{A})$、$(Y,\mathcal{B})$ 为可测空间. 称 $f:X\to Y$ 是可测映射，
若对任意 $B\in\mathcal{B}$ 有 $f^{-1}(B)\in\mathcal{A}$.
\end{definition}

\begin{example}[可测性的有限判别]
设 $\mathcal{B}=\sigma(\mathcal{E})$. 证明 $f$ 可测当且仅当对每个 $E\in\mathcal{E}$
都有 $f^{-1}(E)\in\mathcal{A}$.
\end{example}

\begin{solution}
记 $\mathcal{G}=\{B\subseteq Y:f^{-1}(B)\in\mathcal{A}\}$. 直接验证 $\mathcal{G}$ 是
$Y$ 上的 $\sigma$-代数：它对补与可数并封闭（原像保持），且 $Y\in\mathcal{G}$
因为 $f^{-1}(Y)=X\in\mathcal{A}$. 由假设 $\mathcal{E}\subseteq\mathcal{G}$，
于是 $\mathcal{B}=\sigma(\mathcal{E})\subseteq\mathcal{G}$.
\end{solution}

\begin{remark}
有限判别把「验证无穷多个可测集」压缩成「验证一个生成族」，实际计算时通常取
区间或开集作生成元.
\end{remark}

\section{初等习题}

\begin{exercises}
在可测空间 $(X,\mathcal{A})$ 中判断下列命题是否成立，并给出理由.
\end{exercises}

\begin{exercise}
若 $\mathcal{A}$ 与 $\mathcal{B}$ 都是 $X$ 上的 $\sigma$-代数，则 $\mathcal{A}\cap\mathcal{B}$ 也是.
\end{exercise}
\begin{exercise}
若 $A_n$ 可测，则 $\limsup A_n$ 与 $\liminf A_n$ 均可测.
\end{exercise}
\begin{exercise}
常值函数对任意 $\sigma$-代数 $\mathcal{B}$ 都可测.
\end{exercise}

\begin{analysispath}
$\sigma$-代数之交可以任意多个，$\sigma$-代数之并通常不再是 $\sigma$-代数.
把两个 $\sigma$-代数的并生成的 $\sigma$-代数，是理解独立事件共同信息的第一步.
\end{analysispath}

% 外测度与 Carathéodory 扩张
\section{外测度}

\begin{wenxintishi}[阅读顺序]
外测度是「允许 ∞ 的长度」，只在可数并上满足次可加. 它比测度弱，因而容易构造；
后面再用 Carathéodory 条件把它裁成一个真正的测度.
\end{wenxintishi}

\begin{definition}[外测度]
设 $\mu^{*}:\mathcal{P}(X)\to[0,\infty]$ 满足 $\mu^{*}(\varnothing)=0$ 且对任意
$A\subseteq\bigcup_{n}E_n$ 有 $\mu^{*}(A)\le\sum_n\mu^{*}(E_n)$，则称 $\mu^{*}$
为 $X$ 上的外测度.
\end{definition}

\begin{lemma}[单调性]
外测度单调：若 $A\subseteq B$，则 $\mu^{*}(A)\le\mu^{*}(B)$.
\end{lemma}
\begin{proof}
$B\subseteq B\cup\varnothing\cup\varnothing\cup\cdots$，取 $E_1=B$、其余为空集即可.
\end{proof}

\section{Carathéodory 条件}

\begin{theorem}[Carathéodory 判别]
设 $\mu^{*}$ 为 $X$ 上的外测度，令
\[
\mathcal{M}=\{E\subseteq X:\forall\,A\subseteq X,\ \mu^{*}(A)=\mu^{*}(A\cap E)+\mu^{*}(A\setminus E)\}.
\]
则 $\mathcal{M}$ 是 $\sigma$-代数，且 $\mu^{*}\!|_{\mathcal{M}}$ 是测度.
\end{theorem}

\begin{proposition}[平移不变的外测度]
在 $\mathbb{R}$ 上由区间长度生成的外测度满足 $\mu^{*}(E+t)=\mu^{*}(E)$.
\end{proposition}

\begin{corollary}
Lebesgue 测度完备，且所有 Borel 集都可测.
\end{corollary}

\section{例题与练习}

\begin{example}[构造外测度]
设 $X=\mathbb{N}$，$\mu^{*}(A)=\sum_{n\in A}2^{-n}$. 验证 $\mu^{*}$ 是外测度，
并求出 $\mu^{*}(\mathbb{N})$.
\end{example}
\begin{solution}
空集对应零和；次可加性由非负项重排不变得到. 而
$\mu^{*}(\mathbb{N})=\sum_{n=1}^{\infty}2^{-n}=1$.
\end{solution}

\begin{Exercises}
  \item 证明有限集的外测度为零.
  \item 举出 $\mathbb{R}$ 的一个不可测子集存在的论证思路.
  \item 设 $\mu^{*}(E)<\infty$ 且 $E\in\mathcal{M}$，求 $\mu^{*}(E^{c})$.
\end{Exercises}


\backmatter

\chapter*{术语索引}
\addcontentsline{toc}{chapter}{术语索引}
\vspace{1em}
\begin{multicols}{2}\footnotesize
\begin{description}[leftmargin=1.4em,style=nextline]
  \item[$\sigma$-代数] 对补与可数并封闭的集族.
  \item[可测函数] 原像保持可测性的函数.
  \item[Carathéodory 条件] 用外测度判别集可测的判据.
  \item[Lebesgue 积分] 先简单函数后取极限的积分构造.
  \item[$L^p$ 空间] $p$ 次幂可积函数构成的赋范空间.
\end{description}
\end{multicols}

\BackCover

\end{document}
